% Gerado por scripts/migrate_markdown_to_latex.py.
% Fonte migrada: docs/reunioes/2026-07-31/README.md


\chapter{Reunião de 31 de julho de 2026}



\section{Identificação}


\begin{itemize}
\item \textbf{Tema:} validação da implementação, definição do recorte experimental e
preparação dos benchmarks;
\item \textbf{Participantes identificados na transcrição:} Lucas Galvão e Walter Silva
Oliveira;
\item \textbf{Horário registrado:} início às 16:01 (GMT-03:00);
\item \textbf{Duração da transcrição:} 30 min e 7 s;
\item \textbf{Registro integral compilável:} \href{transcricao-completa.tex}{transcricao-completa.tex};
\item \textbf{Fonte no Google Docs:} \href{https://docs.google.com/document/d/1WBgwj1gDAWh_7VGckHvPhRfSlcRzFIwL2JlnIDhRYbQ/edit?usp=drive_web&tab=t.v1ze8hoqd423}{anotações e transcrição};
\item \textbf{Gravação indicada na ata:} \href{https://drive.google.com/file/d/1-FFyxhkpgkhlesSqULSIaujRSiDnT1UN/view?usp=drive_web}{Google Drive}.
\end{itemize}




A transcrição foi gerada automaticamente e pode conter erros de reconhecimento. A cópia integral foi preservada sem correções; esta página concentra apenas a leitura útil para o projeto.



\section{Leitura executiva}









A reunião definiu uma sequência de trabalho: primeiro demonstrar que a implementação local da atenção reproduz referências conhecidas usando as mesmas entradas, pesos e sementes; depois aplicar uma ou duas técnicas de otimização e compará-las por benchmarks controlados. O recorte recomendado foi o mecanismo de autoatenção, sobretudo suas operações matriciais mais custosas, e não uma Transformer completa. Tabelas, gráficos e texto acadêmico devem ser produzidos junto com os experimentos, em vez de ficarem para o final.



\section{Decisões e direcionamentos}


\begin{enumerate}
\item \textbf{Validar antes de otimizar.} Confrontar a implementação local com cálculo
manual e bibliotecas consolidadas, mantendo entrada, pesos, semente e
tolerância comparáveis.
\item \textbf{Isolar o objeto de estudo.} Priorizar a autoatenção e explicitar que o
projeto não pretende implementar nem otimizar toda a arquitetura Transformer.
\item \textbf{Transformar testes em evidência.} Registrar resultados em tabelas e
gráficos, acompanhados do passo a passo e da fundamentação teórica.
\item \textbf{Comparar poucas técnicas de forma controlada.} Depois do portão de
validação, selecionar uma ou duas otimizações e medi-las contra o mesmo
baseline.
\item \textbf{Começar a redação imediatamente.} Conectar conceitos, método, resultados e
limitações enquanto os experimentos são executados.
\item \textbf{Manter modelos pré-treinados como extensão.} Uma inferência com modelo do
Hugging Face pode enriquecer o trabalho, mas somente se houver tempo depois do
núcleo experimental.
\item \textbf{Confirmar as datas oficiais.} A conversa mencionou novembro como referência
provável para entrega/apresentação, mas sem confirmação institucional.
\end{enumerate}


\section{Reconstrução classificada da orientação}




Esta seção distingue o que aparece na fala do orientador das decisões e interpretações feitas depois. Os intervalos remetem à transcrição integral.



\subsection{Pedidos explícitos do professor}


\begin{itemize}
\item comparar a implementação local com uma biblioteca consolidada, citando
TensorFlow ou Keras, com a mesma entrada e a mesma semente (\texttt{00:00:00-04:12});
\item validar passo a passo antes de aplicar técnicas de otimização (\texttt{00:02:53-05:35});
\item materializar a validação em tabela ou gráfico, comparando cálculo/equação,
implementação manual, PyTorch e TensorFlow (\texttt{00:05:35-08:01});
\item depois da validação, escolher uma ou duas técnicas da metodologia e comparar
os resultados (\texttt{00:08:01-09:18});
\item concentrar o estudo no mecanismo de self-attention e em suas operações de
maior custo, em vez de otimizar uma Transformer inteira (\texttt{00:09:18-14:09});
\item começar a escrever texto corrido imediatamente e fechar as avaliações com
resultados e tabelas (\texttt{00:23:13-24:34}).
\end{itemize}


\subsection{Decisões tomadas durante a reunião}


\begin{itemize}
\item a ordem de trabalho ficou como validação da implementação, seleção das
técnicas, aplicação controlada e comparação;
\item o recorte principal ficou na self-attention, com a limitação de escopo e de
tempo usada como justificativa;
\item se uma camada completa de TensorFlow ocultasse o núcleo, bastaria comparar
uma operação/interação equivalente da atenção (\texttt{00:14:09-15:21});
\item uma demonstração com modelo pré-treinado poderia ser feita somente se
houvesse tempo (\texttt{00:15:21} em diante).
\end{itemize}


\subsection{Dúvidas levantadas na reunião}


\begin{itemize}
\item como garantir que a implementação própria está correta;
\item como definir formalmente o que pode ser chamado de Transformer;
\item até onde é necessário validar a arquitetura, além das operações matemáticas;
\item se TensorFlow expõe o núcleo interno de atenção de forma comparável;
\item se haverá tempo para uma inferência com modelo do Hugging Face;
\item quais são as datas institucionais efetivas de entrega e apresentação.
\end{itemize}


\subsection{Sugestões metodológicas, sem caráter de requisito adicional}


\begin{itemize}
\item usar resultados pequenos e palpáveis para sustentar cada avanço;
\item justificar o recorte pela relevância computacional do mecanismo escolhido e
pelo tempo disponível;
\item usar literatura para selecionar técnicas já conhecidas, sem exigir uma nova
técnica de otimização;
\item escrever durante o desenvolvimento para explicitar como os conceitos se
encadeiam.
\end{itemize}


\subsection{Implementações e validações necessárias derivadas diretamente da reunião}


\begin{itemize}
\item adaptador de layout e execução do mesmo núcleo de atenção em
NumPy/manual, PyTorch e TensorFlow/Keras;
\item casos determinísticos com entradas compartilhadas e tolerância explícita;
\item artefato tabular com resultados e erros;
\item benchmark controlado da self-attention com baseline comum e uma ou duas
técnicas, depois da validação numérica;
\item documentação em prosa de método, resultados e limites.
\end{itemize}


\subsection{Tarefas que podem ficar para depois dos primeiros benchmarks}


\begin{itemize}
\item demonstração de inferência com modelo pré-treinado;
\item ampliação para outros componentes ou para uma arquitetura Transformer
completa;
\item técnicas adicionais além do primeiro conjunto controlado.
\end{itemize}


\subsection{Interpretações posteriores - não são pedidos literais do professor}


\begin{itemize}
\item usar NumPy como referência matemática independente;
\item escolher \texttt{keras.\allowbreak{}ops.\allowbreak{}nn.\allowbreak{}dot\_\allowbreak{}product\_\allowbreak{}attention} como API TensorFlow/Keras;
\item fixar \texttt{float32}, CPU, \texttt{atol=1e-6}, \texttt{rtol=1e-5} e seed \texttt{2026} na validação
cruzada;
\item manter CSV e JSON de metadados versionados;
\item classificar evidências como conceituais, numéricas, algorítmicas ou de
hardware;
\item tratar a RTX 4070 Ti Super, seed \texttt{42}, grade, aquecimento e repetições como
requisitos do plano formal de trabalho, não desta reunião.
\end{itemize}


\section{Situação no repositório em 24/08/2026}



\begin{longtable}{@{}>{\RaggedRight\arraybackslash}p{0.273\linewidth}>{\RaggedRight\arraybackslash}p{0.273\linewidth}>{\RaggedRight\arraybackslash}p{0.273\linewidth}@{}}
\toprule
Encaminhamento da reunião & Situação atual & Evidência ou lacuna \\
\midrule
\endfirsthead
\toprule
Encaminhamento da reunião & Situação atual & Evidência ou lacuna \\
\midrule
\endhead
Validar a matemática da atenção & Implementado no código e nos testes & Há comparação da atenção local com NumPy independente, valores controlados e \texttt{torch.\allowbreak{}nn.\allowbreak{}functional.\allowbreak{}scaled\_\allowbreak{}dot\_\allowbreak{}product\_\allowbreak{}attention} em \href{../../../fases/01_validacao_conceitual/tests/test_validacao_conceitual.py}{test\_\allowbreak{}validacao\_\allowbreak{}conceitual.py}. \\
Comparar com bibliotecas consolidadas & Concluído para o núcleo da atenção & NumPy manual, PyTorch manual, PyTorch SDPA e Keras SDPA com backend TensorFlow foram comparados nos mesmos casos. As 9 comparações passaram; veja a \href{../../../fases/01_validacao_conceitual/evidencias/comparacao_frameworks/comparacao_atencao.tex}{evidência versionada}. \\
Restringir o escopo à autoatenção & Parcialmente alinhado & A validação isola a atenção, mas o executor principal mede um \texttt{Simplified\allowbreak{}Transformer\allowbreak{}Block}; é preciso decidir se o benchmark final será da atenção isolada ou do bloco simplificado. Ver \href{../../../fases/01_validacao_conceitual/README.md}{README da fase}. \\
Preparar benchmarks controlados & Infraestrutura pronta; execução canônica pendente & O executor, o esquema do CSV e os testes de contrato existem em \href{../../../fases/01_validacao_conceitual/validacao/benchmark_controlado.py}{benchmark\_\allowbreak{}controlado.py}, mas não há resultados experimentais versionados. \\
Produzir tabelas, gráficos e conclusões & Automação pronta; artefatos pendentes & \href{../../../fases/01_validacao_conceitual/validacao/analise_resultados.py}{analise\_\allowbreak{}resultados.py} processa os CSVs, mas depende da coleta real. \\
Formalizar a metodologia & Em andamento avançado & A fase já possui protocolo, matriz de ferramentas, perguntas de pesquisa e documentação de validação em \href{../../../fases/01_validacao_conceitual/docs/}{docs da fase}. \\
Usar modelo pré-treinado & Não iniciado e não obrigatório & Deve continuar fora do caminho crítico até que validação, benchmarks e relatório estejam fechados. \\
\bottomrule
\end{longtable}




Essa tabela descreve o que está implementado e documentado; ela não afirma que a bateria final de benchmarks já foi executada no ambiente-alvo.



\section{Próximos passos priorizados}



\subsection{P0 - Fechar o contrato experimental}


\begin{enumerate}
\item Escolher explicitamente o objeto medido no resultado principal: atenção
isolada, como favorecido pela reunião, ou bloco Transformer simplificado.
\item Fixar baseline, técnicas de otimização, shapes, lotes, tamanhos de sequência,
\texttt{dtype}, dispositivo, sementes, aquecimento, repetições e tolerâncias.
\item Preservar TensorFlow/Keras como portão de correção já concluído e reexecutar a
comparação somente quando mudar a atenção, as versões dos frameworks ou as
tolerâncias.
\item Confirmar com a orientação as datas institucionais de submissão, entrega e
apresentação.
\end{enumerate}


\subsection{P1 - Concluir a validação e coletar evidências}


\begin{enumerate}
\item Reexecutar a suíte no ambiente fixado e guardar versão das dependências,
hardware e resultado dos testes.
\item Fazer um piloto curto em CPU para validar o fluxo de ponta a ponta.
\item Executar a grade final no dispositivo-alvo, salvando CSV e metadados. Como
\texttt{resultados/\allowbreak{}} está ignorado pelo Git, definir antes quais artefatos finais
serão versionados ou publicados em outro local reprodutível.
\item Distinguir claramente simulação numérica de ganho efetivo de hardware; só usar
esse último rótulo quando armazenamento, tipo numérico e kernel executado
sustentarem a afirmação.
\end{enumerate}


\subsection{P1 - Escrever o relatório junto com a análise}


\begin{enumerate}
\item Gerar uma tabela de validação com entrada, saída esperada, saída local, saída
da referência, erro máximo, tolerância e status.
\item Gerar tabelas e gráficos de latência, throughput, memória, erro e redução de
operações para cada cenário contra o mesmo baseline.
\item Responder às perguntas de pesquisa com os resultados observados e registrar
ameaças à validade, limitações e casos em que uma técnica não trouxe ganho.
\item Incorporar método, resultados e discussão ao texto acadêmico sem esperar o
encerramento de toda a coleta.
\end{enumerate}


\subsection{P2 - Extensão opcional}





Depois de concluir o núcleo, avaliar uma demonstração de inferência com um modelo pré-treinado do Hugging Face. A extensão não deve alterar o baseline nem atrasar as entregas principais.



\section{Critério de conclusão desta etapa}



A etapa pode ser considerada fechada quando houver:


\begin{itemize}
\item validação reproduzível da atenção contra as referências escolhidas;
\item decisão documentada sobre o escopo e sobre TensorFlow/Keras;
\item benchmarks comparáveis, com configuração e ambiente registrados;
\item CSVs, tabelas e gráficos rastreáveis até a execução;
\item conclusões compatíveis com o nível real de evidência, sem extrapolação para
hardware ou Transformer completa;
\item seção correspondente já integrada ao relatório do projeto.
\end{itemize}


\section{Trechos da reunião que sustentam os encaminhamentos}


\begin{itemize}
\item \textbf{00:00-04:12:} necessidade de validar a implementação contra bibliotecas
consolidadas antes de avançar;
\item \textbf{05:35-08:01:} sugestão de tabelas e gráficos comparando cálculo manual,
PyTorch e TensorFlow, além de iniciar o relatório;
\item \textbf{08:01-14:09:} sequência validação → otimização → comparação e delimitação do
escopo na autoatenção;
\item \textbf{15:21-16:15:} modelo pré-treinado como possibilidade condicionada ao tempo e
incerteza sobre o calendário de novembro;
\item \textbf{23:13 em diante:} recomendação de redigir em prosa desde já e fechar cada
avaliação com resultados e tabelas;
\item \textbf{29:30-30:07:} compromisso de avançar nos benchmarks após a reunião.
\end{itemize}
